\chapter{Въведение}

\section{Случаен експеримент и~елементарни събития}

\begin{defbox}[~(Случаен експеримент)]
\textbf{Случаен експеримент} е опит или наблюдение, чийто изход не може
да бъде предсказан предварително. Всеки възможен елементарен резултат
от този експеримент наричаме \textbf{елементарно събитие} и~го
означаваме с~малка гръцка буква (обикновено~$\omega$).
\end{defbox}

\begin{rembox}[Мотивация]
Случайният експеримент е моделът, чрез който математиката хваща
явления с~недетерминиран изход - от хвърляне на монета и~зар, през
резултати от избори, до напълно непрекъснати процеси. Класически
пример за такъв процес е \textbf{Брауновото движение}: през 1827~г.
ботаникът Робърт Браун наблюдава хаотичното движение на прашец във
вода, а~през 1905~г. Айнщайн му дава строго вероятностно обяснение.
Именно нуждата от такива по-сложни пространства ни води до общата
дефиниция по-долу.
\end{rembox}

\begin{defbox}[~(Пространство от елементарни събития)]
\textbf{Пространството от елементарни събития} е множеството от
всички възможни елементарни събития за даден случаен експеримент.
Означава се с~главна гръцка буква (обикновено~$\Omega$).
\end{defbox}

\begin{defbox}[~(Събитие)]
Нека $\Omega$ е пространството от елементарни събития за даден
случаен експеримент. Всяко подмножество $A \subseteq \Omega$ наричаме
\textbf{събитие}. Едно събитие \emph{настъпва}, ако резултатът
$\omega$ от експеримента е елемент на това подмножество, т.е.
$\omega \in A$. Най-често събитието групира елементарни събития по
даден признак (рядко е точно едно елементарно събитие).
\end{defbox}

\begin{rembox}[Операции с~events (кратко резюме)]
Събитията са множества, затова върху тях работят обичайните
операции и~закони от теория на множествата:
\begin{itemize}[nosep]
  \item $A \subseteq B \iff \forall \omega \in A \Rightarrow \omega \in B$ (влагане); $A = B \iff A\subseteq B \text{ и~} B \subseteq A$.
  \item $A \cup B = \{\omega \in \Omega : \omega \in A \text{ или } \omega \in B\}$ - \emph{„или“};
  \item $A \cap B = \{\omega \in \Omega : \omega \in A \text{ и~} \omega \in B\}$ - \emph{„и~“};
  \item $A^c = \{\omega \in \Omega : \omega \notin A\}$ - \emph{допълнение} (по-рядко: отрицание), като $\Omega = A \cup A^c$;
  \item $A \setminus B = A \cap B^c$ (,,$A$ без $B$``);
  \item комутативност: $A\cup B = B\cup A$, $A\cap B = B\cap A$;
  \item асоциативност: $A\cup(B\cup C) = (A\cup B)\cup C$;
  \item дистрибутивност: $A\cup(B\cap C) = (A\cup B)\cap(A\cup C)$ и~двойственото ѝ равенство;
  \item инволутивност: $(A^c)^c = A$;
  \item закони на Де Морган: $\left(\bigcup_{i=1}^{\infty} A_i\right)^c = \bigcap_{i=1}^{\infty} A_i^c$, \quad $\left(\bigcap_{i=1}^{\infty} A_i\right)^c = \bigcup_{i=1}^{\infty} A_i^c$.
\end{itemize}
\end{rembox}

\section{\texorpdfstring{$\sigma$}{Sigma}-алгебра}

\begin{defbox}[~(\texorpdfstring{$\sigma$}{Sigma}-алгебра)]
Нека $\Omega$ е произволно множество (наричано пространство от
елементарни събития) и~нека $\mathcal{A}$ е система от подмножества
на $\Omega$ (т.е. $\mathcal{A} \subseteq \mathcal{P}(\Omega)$).
Казваме, че $\mathcal{A}$ е \textbf{$\sigma$-алгебра} на $\Omega$, ако
са изпълнени:
\begin{enumerate}[label=\arabic*)]
  \item $\Omega \in \mathcal{A}$ ($\mathcal{A}$ съдържа цялото пространство);
  \item ако $A \in \mathcal{A}$, то $A^c = \Omega \setminus A \in \mathcal{A}$ (затвореност относно допълнение);
  \item ако $A_1, A_2, \dots \in \mathcal{A}$, то $\bigcup_{i=1}^{\infty} A_i \in \mathcal{A}$ (затвореност относно изброимо обединение).
\end{enumerate}
\end{defbox}

\begin{corbox}[~(основни свойства)]
Ако $\mathcal{A}$ е $\sigma$-алгебра на $\Omega$, то:
\begin{itemize}[nosep]
  \item $\varnothing \in \mathcal{A}$;
  \item ако $A_i \in \mathcal{A}$ за всяко $i \geq 1$, то $\bigcap_{i=1}^{\infty} A_i \in \mathcal{A}$ (затвореност и~относно изброимо сечение);
  \item ако $A_1,\dots,A_n \in \mathcal{A}$, то $\bigcup_{i=1}^{n} A_i \in \mathcal{A}$ и~$\bigcap_{i=1}^n A_i \in \mathcal A$ (затвореност и~относно \emph{крайно} обединение/сечение).
\end{itemize}
\end{corbox}

\begin{proofbox}
Понеже $\Omega \in \mathcal{A}$ и~$\Omega^c = \varnothing$, от аксиома~2)
следва $\varnothing = \Omega^c \in \mathcal{A}$.

Нека $A_i \in \mathcal{A}$ за всяко $i \geq 1$. От аксиома~2) $A_i^c \in
\mathcal{A}$ за всяко $i$, а~от аксиома~3), приложена към редицата
$\{A_i^c\}$, следва $\bigcup_{i=1}^{\infty} A_i^c \in \mathcal{A}$.
Прилагайки отново аксиома~2) към това обединение и~закона на Де
Морган:
\[
  \left(\bigcup_{i=1}^{\infty} A_i^c\right)^{c} \in \mathcal{A}
  \quad\Longrightarrow\quad
  \bigcap_{i=1}^{\infty} A_i \in \mathcal{A}.
\]

За крайното обединение: нека $A_1,\dots,A_n \in \mathcal{A}$. Допълваме
редицата с~$A_{n+k} := \varnothing$ за $k \geq 1$ (знаем вече, че
$\varnothing \in \mathcal{A}$). Тогава по аксиома~3)
\[
  \bigcup_{i=1}^{n} A_i = \bigcup_{i=1}^{\infty} A_i \in \mathcal{A}.
\]
Затвореността относно крайно сечение следва по същия начин чрез
допълнение и~закона на Де Морган. \qedbg
\end{proofbox}

\begin{rembox}[Примери]
Нека $\Omega = \{a,b\}$. Най-простата възможна $\sigma$-алгебра е
\textbf{тривиалната} $\mathcal{A}_1 = \{\varnothing,\Omega\}$. От
другата страна, \textbf{дискретната} (пълна) $\sigma$-алгебра е
$2^{\Omega} = \{\varnothing,\Omega,\{a\},\{b\}\}$ - множеството от
\emph{всички} подмножества на $\Omega$; тя винаги е $\sigma$-алгебра
и~съдържа $2^{|\Omega|}$ елемента. За произволна крайна извадка от $n$
изхода имаме $|2^\Omega| = 2^n$ възможни събития.
\end{rembox}

\section{Топология и~Борелова \texorpdfstring{$\sigma$}{sigma}-алгебра}

\begin{defbox}[~(Система от отворени множества)]
Нека $X$ е произволно множество. Система от отворени множества
(\textbf{топология}) е съвкупност $\tau \subseteq 2^X$, която
удовлетворява аксиомите:
\begin{enumerate}[label=\arabic*)]
  \item Празното множество и~цялото пространство са отворени.
  \item Обединението на произволен брой отворени множества е отворено.
  \item Сечението на \emph{крайно много} отворени множества е отворено.
\end{enumerate}
\end{defbox}

\begin{defbox}[~(Породена / Борелова \texorpdfstring{$\sigma$}{sigma}-алгебра)]
Нека $\Omega$ е множество от елементарни събития и~$G$ е произволна
колекция от подмножества на $\Omega$. Най-малката $\sigma$-алгебра,
съдържаща $G$, се нарича \textbf{породена от $G$ $\sigma$-алгебра} и~се означава с~$\sigma(G)$:
\[
  \sigma(G) := \bigcap_{\substack{\mathcal{A} \supseteq G\\ \mathcal{A}\ \sigma\text{-алгебра}}} \mathcal{A}.
\]
Ако $(X,\tau)$ е топологично пространство, \textbf{Бореловата
$\sigma$-алгебра} е $\mathcal{B}(X) := \sigma(\tau)$ - най-малката
$\sigma$-алгебра, съдържаща всички отворени множества на~$\tau$.
\end{defbox}

\begin{rembox}[Пример: породена \texorpdfstring{$\sigma$}{sigma}-алгебра от три събития]
Нека $A,B,C$ образуват пълна група от несъвместими събития в~$\Omega$ (т.е.\ $A,B,C$ са по двойки непресичащи се и~$A\cup B\cup
C=\Omega$) и~$G=\{A,B,C\}$. Тогава породената $\sigma$-алгебра е
крайна:
\[
  \sigma(G) = \{\varnothing,\Omega,A,B,C,A\cup B,A\cup C,B\cup C\}.
\]
Така изброимо-безкрайният проблем се редуцира до събиране на краен
брой непресичащи се „тухлички“.
\end{rembox}

\begin{rembox}[Пример: Бореловата \texorpdfstring{$\sigma$}{sigma}-алгебра на \texorpdfstring{$\mathbb{R}$}{R}]
За $\Omega = \mathbb{R}$ с~обичайната топология, $\mathcal{B}(\mathbb{R})$
съдържа всички интервали $(a,b)$, $(a,\infty)$, $(-\infty,a)$ за
$a,b\in\mathbb{R}$, $a<b$. Съдържа и~всяко едноточково множество
$\{a\}$, защото
\[
  \{a\} = \bigcap_{i=1}^{\infty} \left(a-\tfrac1i,\, a+\tfrac1i\right) \in \mathcal{B}(\mathbb{R}),
\]
а~оттам и~$[a,b) = \{a\}\cup(a,b) \in \mathcal{B}(\mathbb{R})$,
$[a,b]=\{a\}\cup(a,b)\cup\{b\}\in\mathcal B(\mathbb R)$.
\end{rembox}
